\documentclass[twocolumn,10pt]{article}
\usepackage[utf8]{inputenc}
\usepackage[margin=1in]{geometry}
\usepackage{graphicx}
\usepackage{amsmath}
\usepackage{hyperref}
\usepackage{xcolor}
\usepackage{booktabs}
\usepackage{cite}

\title{RainbowZero: A Unified Framework for Planning and Learning in Stochastic and Multi-agent Environments}
\author{Research Document}
\date{\today}

\begin{document}

\maketitle

\begin{abstract}
Monte Carlo Tree Search (MCTS) combined with deep neural networks has achieved superhuman performance in perfect-information games, yet extending these methods to stochastic and multi-agent environments remains a significant challenge. We present RainbowZero, a modular and extensible framework that unifies MuZero and its most significant enhancements, including Stochastic MuZero for aleatoric uncertainty, Gumbel MuZero for deterministic policy improvement, and EfficientZero for improved sample efficiency. Our framework introduces novel extensions to accommodate $N$-player dynamics and irregular turn structures, enabling planning in complex environments previously beyond the scope of standard MuZero. We evaluate RainbowZero across a diverse suite of benchmarks, ranging from classic control tasks like CartPole to the stochastic multi-agent board game Catan. Our results demonstrate that the integrated architectural improvements significantly enhance sample efficiency and robustness to stochasticity, providing a state-of-the-art baseline for general-purpose model-based reinforcement learning.
\end{abstract}

\section{Introduction}
\begin{itemize}
    \item Motivation for a unified MuZero framework.
    \item Challenges in reinforcement learning: stochasticity, multi-agent dynamics, and sample efficiency.
    \item Summary of contributions: unified implementation, support for $>2$ players, and integration of EfficientZero, Gumbel MuZero, and Stochastic MuZero.
\end{itemize}

\section{Related Work}
\begin{itemize}
    \item \textbf{MuZero}: The foundation of model-based RL without a known simulator [CIT1].
    \item \textbf{Stochastic MuZero}: Handling aleatoric uncertainty via afterstates and chance nodes [CIT2].
    \item \textbf{Gumbel MuZero}: Policy improvement with few simulations and deterministic guarantees [CIT3].
    \item \textbf{EfficientZero \& V2}: Sample efficiency through consistency losses and value prefixes [CIT4, CIT5].
    \item \textbf{Sampled MuZero}: Scaling to large or continuous action spaces [CIT6].
    \item \textbf{MuZero Unplugged}: Learning from offline datasets and reanalysis [CIT7].
\end{itemize}

\section{Methodology}
RainbowZero integrates several advancements into a single, modular codebase.

\subsection{Base MuZero Architecture}
\begin{itemize}
    \item Representation, Dynamics, and Prediction networks.
    \item Generalized MCTS for planning.
\end{itemize}

\subsection{Stochasticity and Afterstates}
\begin{itemize}
    \item Implementation of \textbf{Stochastic MuZero}.
    \item Introduction of \textbf{Chance Nodes} in the search tree.
    \item Learning afterstates to decouple transition dynamics from stochastic outcomes.
\end{itemize}

\subsection{Enhanced Search Strategies}
\begin{itemize}
    \item \textbf{Gumbel MuZero}: Integration of Gumbel-Top-k sampling and Sequential Halving.
    \item \textbf{Sampled MuZero}: Handling large action spaces by sampling a subset of actions for expansion.
\end{itemize}

\subsection{Sample Efficiency (EfficientZero Integration)}
\begin{itemize}
    \item \textbf{Value Prefix}: Using LSTM-based reward prediction for better long-term credit assignment.
    \item \textbf{Self-Consistency}: A SSL-style loss to ensure latent state consistency across dynamics steps.
    \item \textbf{Distributional RL}: Support-based value and reward predictions.
\end{itemize}

\subsection{Unified Multi-Agent Support}
\begin{itemize}
    \item \textbf{Beyond 2 Players}: Support for N-player games (e.g., Catan).
    \item \textbf{Non-Alternating Turns}: Handling environments where the same player may move multiple times or turns are irregular.
    \item \textbf{Player Wrappers}: Unified interface for multi-agent environments.
\end{itemize}

\section{Experimental Setup}
\subsection{Environments}
\begin{itemize}
    \item \textbf{CartPole}: Classic control benchmark for basic convergence.
    \item \textbf{Tic-Tac-Toe}: Simple perfect-information game for verifying MCTS and multi-player logic.
    \item \textbf{Slippery Grid}: A variation of the Cliff Walk environment designed to test stochastic modeling.
    \item \textbf{Catan}: A complex, multi-agent board game with stochasticity (dice rolls) and large action spaces (trading, building).
\end{itemize}

\subsection{Implementation Details}
\begin{itemize}
    \item Hardware: [Specify GPU/CPU used].
    \item Hyperparameters: [Reference Appendix].
\end{itemize}

\section{Results and Analysis}

\subsection{Baseline Comparisons on CartPole}
Comparison of MuZero, EfficientZero, and Stochastic MuZero.
\begin{figure}[h]
    \centering
    \fbox{Graph Placeholder: CartPole Training Curves}
    \caption{Training curves for all models on CartPole. Expectation: EfficientZero should show faster initial convergence.}
    \label{fig:cartpole}
\end{figure}

\subsection{Strategic Excellence in Tic-Tac-Toe}
Model playing against expert/random agents.
\begin{figure}[h]
    \centering
    \fbox{Graph Placeholder: Tic-Tac-Toe expert Win Rate}
    \caption{Win rate vs Expert Agent in Tic-Tac-Toe.}
    \label{fig:tictactoe}
\end{figure}

\subsection{Modeling Uncertainty in Slippery Grid}
Comparison of Stochastic vs. Non-stochastic models.
\begin{figure}[h]
    \centering
    \fbox{Graph Placeholder: Slippery Grid expected Reward}
    \caption{Expected Reward vs. Stochasticity Level in Slippery Grid. Expectation: Stochastic MuZero should significantly outperform deterministic variants as the "slip" probability increases.}
    \label{fig:grid}
\end{figure}

\subsection{Scaling to Catan}
Performance in a 3-4 player setting.
\begin{figure}[h]
    \centering
    \fbox{Graph Placeholder: Catan Win rates}
    \caption{Win rates and resource collection efficiency in Catan.}
    \label{fig:catan}
\end{figure}
Analysis of turn-taking and stochastic dice roll modeling.

\section{Discussion}
\begin{itemize}
    \item Trade-offs between model complexity and performance.
    \item Effectiveness of the unified framework across different game types.
    \item Insights into multi-agent dynamics in Catan.
\end{itemize}

\section{Conclusion}
\begin{itemize}
    \item Summary of findings.
    \item Future work: Continuous control, real-world robotics applications.
\end{itemize}

\section{References}
[Citations to be added here]

\end{document}
